\section{Limitations, investigated}

The single-mode limitations of the engine stand unchanged: a \(\chi\) read from a lone real pole with no partner is undefined (\(\omega_0\) is not recoverable from one pole) and is reported as such rather than invented; mode parentage must be read together with the frequency gap, since a mid-range parentage with an open gap indicates hybridization while a closing gap indicates coalescence, and parentage near one-half does not by itself certify an exceptional point; non-reciprocal (asymmetric \(C\) or \(K\)) systems are rejected at the input contract rather than returned as spurious mechanical modes; and negative damping, non-Hermitian mass, and many-mode (more than three) coupling are untested, with no robustness claimed.

The extension of Part~III adds its own bounded items, stated rather than hidden. Step~1's resolvability threshold is line-shape- and criterion-specific: the equal-Gaussian merger near $\Delta \approx 2\sigma$ is illustrative, not a universal constant, and two unresolved lines are unresolved under a stated line shape and criterion rather than a mathematical continuum. Step~2's $R_\phi = 1$ is the equal-contribution crossover, not a critical point and not an identifiability transition, and is claimed as such: with $T_1$ measured, $\chi$ is recoverable at every $R_\phi$, and $R_\phi$ reports the bias ($1 + R_\phi + R_{\mathrm{or}}$) of the naive width reading rather than a boundary of recoverability. The reversibility of coherence loss is not binary: an echo refocuses static inhomogeneity completely, fast stochastic dephasing largely not, and intermediate correlation times partly, so reversibility diagnoses correlation time and does not by itself define an axis. The single worked example (Rh(CO)$_2$acac) is anchored to the two endpoints the source reports in text (3.4~K: $T_2 = 95.2$~ps, width 0.11~cm$^{-1}$; 250~K: width 1.5~cm$^{-1}$ (source text), from which $T_2 = 7.08$~ps follows by Eq.~1 and is author-derived rather than stated). The intermediate echo and pump-probe points are by-eye transcriptions of the source's Figure~2 and are shown as an illustrative trajectory only; no crossover temperature or intermediate bias factor is reported as a result. The quantitative claim is limited to the endpoint contrast, a homogeneous-width growth of about 13.6; the fitted-lifetime $\chi$ near $3\times10^{-5}$ at both endpoints is source-informed illustrative rather than source-anchored, pending calibrated pixel extraction of the pump-probe points, and the mechanical-placement claim is scoped to this single example with no coupled-pair catalogue claimed. A calibrated redigitization of Figure~2 would be required before the lifetime-derived $\chi$ or bias could be reported quantitatively. The plasma discussion is an analogy, not a worked entry: the free-streaming calculation illustrates reversible phase mixing but is not a self-consistent Vlasov--Poisson solution and does not reproduce the Landau pole.

A separate boundary applies when the generator is not supplied but inferred. Subspace identification from a vector response record recovers the spectrum accurately but does not by itself certify the parent's multiplicity structure: generic finite-data perturbation splits a defective pole, so a fitted generator is a reduced representation in the sense of Section~2.3 and its structural verdict is scoped to the fit. No such estimator is shipped here.

For general first-order nonnormal generators, pole classification does not determine finite-time transient amplification or pseudospectral sensitivity. ChemSA presently reports spectral structure and does not claim a nonmodal stability bound; a pseudospectral or propagator-norm extension remains future work.
